\section{Estado del Arte}
\label{sec:estado_del_arte}

Este capítulo analiza el panorama actual de los modelos de lenguaje, la evolución hacia modelos más eficientes y las investigaciones previas en arquitecturas híbridas.

\subsection{Evolución de los Modelos de Lenguaje: De LLM a SLM}
    En los últimos años, el campo del Procesamiento del Lenguaje Natural (NLP) ha estado dominado por los \textit{Large Language Models} (LLM), cuya arquitectura base sigue siendo el Transformer propuesto por Vaswani et al. \cite{vaswani2017attention}. Modelos como Llama 2 \cite{touvron2023llama} han demostrado capacidades emergentes sorprendentes al escalar el número de parámetros a miles de millones. 

Sin embargo, recientemente ha surgido una tendencia hacia los \textit{Small Language Models} (SLM), impulsada por la necesidad de eficiencia y despliegue en entornos con recursos limitados. Investigaciones como \textit{"Textbooks Are All You Need"} \cite{gunasekar2023textbooks} demuestran que la calidad de los datos puede compensar el tamaño del modelo, permitiendo que arquitecturas de menos de 3 mil millones de parámetros (como Phi-3 o Mistral) compitan en tareas específicas con modelos mucho más grandes. En este contexto, la familia de modelos SmolLM \cite{allal2024smollm} representa el estado actual de la técnica en SLMs ultra-eficientes, optimizando tanto el proceso de entrenamiento como la curación de datos para maximizar el rendimiento en dispositivos de consumo.

\subsection{Limitaciones de los Transformers Puros}
A pesar de su éxito, los Transformers convencionales presentan un cuello de botella fundamental: la complejidad cuadrática del mecanismo de auto-atención respecto a la longitud de la secuencia \cite{vaswani2017attention}. Esto implica un consumo de memoria de video (VRAM) y un tiempo de cómputo que crece de forma prohibitiva para contextos largos. En un entorno de producción, esto se traduce en altos costes operativos y latencias elevadas, lo que motiva la búsqueda de arquitecturas con escalado lineal.

\subsection{Arquitecturas Híbridas y Recurrencia Moderna}
Para mitigar las limitaciones del Transformer, han resurgido las arquitecturas basadas en recurrencia, pero con enfoques modernos. Las \textit{Gated Recurrent Units} (GRU) \cite{cho2014learning} ofrecen un manejo eficiente de secuencias con un estado oculto de tamaño fijo, eliminando el coste cuadrático. 

Investigaciones recientes como Mamba \cite{gu2023mamba} han propuesto modelos de espacio de estados (\textit{State Space Models}) que logran un rendimiento comparable a los Transformers pero con inferencia de tiempo lineal. La tendencia actual, y el núcleo de este TFM, es la hibridación: combinar capas de atención para capturar relaciones globales con capas recurrentes o lineales para mantener la eficiencia secuencial.

Esta hibridación se ha visto facilitada por la evolución de los ecosistemas de software, específicamente con la transición de la librería \textit{Transformers} de Hugging Face hacia la versión 5. Mientras que las versiones precedentes (v4.x) requerían implementaciones manuales complejas para integrar componentes recurrentes en el grafo de computación del Transformer, la versión 5 introduce abstracciones nativas para modelos de espacio de estados y arquitecturas secuenciales no atencionales. Estas mejoras permiten una gestión optimizada de la caché de inferencia (KV Cache) y una integración transparente de capas GRU dentro del flujo autorregresivo, garantizando que el modelo mantenga la compatibilidad con herramientas de optimización y despliegue estándar del sector.

\subsection{Aumento por Recuperación (RAG)}
Finalmente, el estado del arte no solo se centra en la arquitectura del modelo, sino en su interacción con fuentes externas. El \textit{Retrieval-Augmented Generation} (RAG) \cite{lewis2020rag} se ha consolidado como el estándar para reducir alucinaciones en modelos de lenguaje. Esta técnica permite que el modelo consulte bases de datos externas en tiempo real, proporcionando respuestas basadas en evidencia sin necesidad de reentrenamientos costosos, un componente vital para asistentes técnicos especializados.
